\section{Interview with Aphrodite}

There is spiritual current today that often goes under the name ``The Law of Attraction'', although by other names in the past. This was brought to prominence in a documentary made a few years ago called the ``Secret''. Its main spokesmen include Oprah Winfrey, Deepak Chopra, Wayne Dyer, the author of ``Conversations with God'', and so on. While employing, without understanding, some legitimate Hermetic practices, the overall thrust of this movement is false.

Rather than the rejection of materiality as seen, for example, in Theosophy, this movement seeks to improve the material conditions of life through some spiritual practices which include affirmations and visualizations. My son was invited to a Wayne Dyer talk several years ago; he told me the audience included almost entirely post-menopausal women. These are woman who have achieved a measure of success in the world and regard it as a sign of superior spirituality.

\textbf{Julius Evola} described several ``races of the spirit'' thatcoften display considerable insight into the human condition. In his description of the aphroditic race, he pinpoints the relevant features that describe this type. There is no sense of transcendence, so the focus is on improving material existence: money, relationships, property, etc. The end result is the confusion of intense sensual experiences with spiritual realization. This is Evola's description:

\begin{quotex}
Another race of the spirit is the truly aphroditic race; in it, tellurism, i.e., the adherence to terrestriality, assumes the forms of an extreme refinement of material existence, and not seldom goes on to promote an opulent development of everything that is pomp and luxury in the outer life, therefore also of the world of arts and aesthetic sentiment. But, in the inner life, a passivity and a lunar inconsistency subsist, balanced by a particular prominence given to sensuality but also to all that which is related to woman, which, through such a way, goes on to exercise control and to silently secure for itself a superior position.
\end{quotex}


\paragraph{The Massage Queen}


I recently had the pleasure of an hour long interview with an exemplar of the aphroditic race. In a narrow sense, she can be said to have won the lottery of life. In her 50s, she was still attractive, healthy, financially independent, extroverted, with the spare time to pursue her numerous hobbies. Knowing my interests, she was eager to share our spiritual perspectives.

Opening up the topic, she surprised me when she said that the spiritual turning point in her life came from reading the book \emph{Think and Grow Rich} by \textbf{Napoleon Hill}. This motivated her to open up some spas in the Midwest where she was able to combine her interest in, and knowledge of, the chakras with massage. She eventually sold them at a tidy profit and retired to Florida.

Now such a book is not a spiritual book at all, despite its distortions of certain Hermetic exercises. Rather, it acts at the psychological level by motivating someone to single-pointedly focus on making money. Furthermore, for anyone who is convinced by it and harbors no doubts, it can bring a lot of confidence to that project. She was such a person, and the results are there to see. Nevertheless, Hill's book is not a passive teaching, of the sort that the ``universe will attract it to you'', but a call to action.

That led to a discussion of the so-called Law of Attraction, which is a generalization of Hill's book. Not just money, but health, love relationships, and all other aspects of life can be materialized in the same way. Eventually, however, those of that mindset tire of the effort required, and decide that one just needs to recite affirmations and results will follow effortlessly. Of course, work and toil are part of our lot and cannot be dispensed with. Now I am not discouraging anyone from making affirmations, which have their use provided they lead to manifestation through the will and the application of energy.

Keep in mind, however, that the Law of Attraction cannot attract sanctity to you, which is really what should concern you. At one point, she suggested that I attend a nearby Buddhist center. I didn't tell here that I have been initiated into a Tibetan Buddhist lineage, but I did bring up the point of a common meditation theme. In that, one focuses on a mandala and envisions oneself in the company of Buddhas, Bodhisattvas, gods, and other higher beings. I explained that we unconsciously throughout the day create a lower self-image; this meditation breaks that chain by identifying the self with higher beings. She gave me a slight frown and changed the topic, as though she was thinking, ``why would I want to give up my self?''

My point is to illustrate the aphroditic race. This woman's spirituality was focused on ``the extreme refinement of material existence'' and ``an opulent development of everything that is pomp and luxury in the outer life''.

Somehow, out of nowhere, there is a new word, ``bossy'', that politically correct people should not use today. That is fine, because what they are trying to define out of existence, or ignore, is what Jung described as the ``animus''. This is the unconscious masculine element in a woman's psyche that can emerge in a way that can be sometimes be described as ``bossy''. I had to navigate around this woman's animus in order to keep the conversation light and pleasant.

She was always the ``first'' to do something, the innovator, etc. While she should be proud of her accomplishments in bringing hot stones to massage therapy, she did not discover chakras. Moreover, a man of Tradition would be horrified to discover he was an ``innovator'', since his aim is to extend Tradition. As she posed questions to me, no matter how I answered, she would assert that she ``already knew that''.

Obviously, I could tell what she understood and what she did not. I also know how difficult it is to achieve any worthwhile insights on the spiritual path, and pop psychology is a bane, not a boon. When I tried to describe Guenon's ideas on tradition, she immediately assumed I was referring to syncretism. She pointed out that not only did she have chakra charts and other Oriental symbols decorating her home, she also had a crucifix.

Not surprisingly, the conversation turned to relationships. Having in mind the Mouravieff's teaching on polar beings, I asked her if she believed that a man and a woman who are spiritually united could achieve more together than apart. Her immediate, but irrelevant response, abruptly ended the conversation.

She told me yes. She related how she once achieved such an intense orgasm that it raised up her kundalini all the way to the top. And the irony, she explained, is that the guy involved knew nothing at all about metaphysics, chakras, or kundalini!

That is what it comes down to. Despite her exalted self-image and spiritual path, what she was really looking for was a good f@\#k. She could only conceive of kundalini in material terms. Truth be told, I have run into this attitude in several such women. Evola describes the telluric race in these terms:

\begin{quotex}
Sexuality has a considerable part in it, in its most elementary aspect: naturally, not just phallic, virile sexuality; in this regard, it can actually be said that it can rather more easily appear in a woman than in a man, in accordance with an entirely telluric nature.
\end{quotex}



\flrightit{Posted on 2014-03-17 by Cologero}

\begin{center}* * *\end{center}

\begin{footnotesize}
\begin{sffamily}
\texttt{Matt on 2014-03-17 at 02:24 said:}

There is a Protestant equivalent of this among pentecostals/charismatics and other evangelicals. Notions of ``God's bank account'' and ``God wants you to be rich and successful'' are prevalent; this probably shouldn't be confused with the Calvinist idea of wealth as the sign of the elect as the two currents seem to have a different focus (improvement of material conditions and opulence for the former, concern for the celestial state in the case of the latter). I don't know what race would be linked to that Calvinist thought.

Your conversation with that women reminds me of a similar conversation I had with my sister over dinner a year ago. She told me she was into ``meditation'' and yoga (what woman in their 20's living in the west isn't into that?). I then asked her the obvious question as to why she engages in it. Her response was that it improves one's physical and mental health, and then she cited some various studies that suggested evidence for that. I told her while that all might be true, the purpose of those exercises is to try and connect one back to a person's transcendent root in the Divine, and those pleasent sides effects are just that -- side effects. To do all of that just for those side effects is along the lines of someone taking prescprition pills just for the possible pleasent side effects of the drug, rather than taking them for their intended, proscribed use. Suffice to say, she got defensive and aggravated pretty quickly. Maybe she just took offense at being related to a prescription drug abuser, or tying this to the point of the post, something more deeper may have been going on. I decided to change the subject of the conversation as there was no point in continuing that specific exchange.

I don't know if she still involves herself in any of that. I find myself becoming estranged from her as the years go on.

\hfill

\texttt{Cologero on 2014-03-17 at 08:21 said:}

Matt, I'm sorry to hear about your estrangement from your sister and hope the rift can be healed. The point of this knowledge (gnosis) is not to use it as a weapon for an ``attack'', but rather to understand that there are hard limits to the possibilities of understanding for certain spiritual types. As we learn to recognize the various types, we are able to navigate the world more successfully.

But I know what you mean about the girls of yoga who want to look hot in their yoga suits ... and they do. You could point them to Patanjali if you like, but there is a curious blindness in the human condition that will hide the plain meaning of the text.

Because we know and understand these limits, we have never been involved in ``debates'' or attempts to convert anyone. Gnosis is certain only for those who have made the requisite efforts, and no one who has not done so can be convinced by empirical evidence or logical proofs. The ultimate task is to be able to recognize those who do have the possibility to manifest higher states. The question of why some are drawn to this work while most are not is a bit of a mystery, and may be worth pondering. But for those who are, it would be a shame to squander the opportunity and only go part of the way.

I now see that it will be useful to go through Evola's descriptions of the various spiritual races, so I will be publishing that material over the next few weeks. Of course, it is even more helpful to illustrate them with specific examples.

\hfill

\texttt{scardanelli on 2014-03-17 at 09:11 said:}

Cologero,

Does this understanding of the spiritual races have anything to do with the so called white race, red race, black race, etc, often discussed by spiritualists/anthroposophists? I haven't given much serious thought to it in the past, but I've come across it in other authors who I trust somewhat more- Mouni Sadhu in Tarot and Sedir in his piece on the history of the Rose Croix (I think).

\hfill

\texttt{UnPoeteMaudit on 2014-03-17 at 10:55 said:}

To go off on a tangent, even as far as the mere realisation of a desire into the physical world is concerned, the `Law of Attraction' teaches a half-baked method that flies in the face of elementary (and tried and tested to practicing occultists at least) principles of practical magic, which insist that the ritual process end in complete detachment from the desire in order for it to manifest.

\hfill

\texttt{Cologero on 2014-03-17 at 12:32 said:}

That's not a tangent, Poet, that is a good observation. ... should have brought you along.

\hfill

\texttt{Cologero on 2014-03-17 at 13:12 said:}

Scardanelli, Evola did not regard biological race as normative. Thus a white-skinned European is not necessarily a spiritual Aryan. It is not a question of who you can trust, but rather of what you can verify. Perhaps it is best to wait until I publish more of his descriptions and we'll see what results from that. It seems to me that his description of the Aphroditic type is pretty accurate, even if the woman in question is of the ``white race''.

\hfill

\texttt{Dan on 2014-12-24 at 08:01 said:}

I really enjoyed reading this article, it provided a down-to-earth image of one of the spiritual races Evola described, by providing modern examples. The way that you and Evola describe the Aphroditic race reminds me of the modern neopagan movements like Wicca. They aren't really focused on any true element of transcendence, solely on refining their material life.

\hfill

\texttt{Alistair Fraser on 2019-03-17 at 06:18 said:}

Take a moment to indulge in some dark humour this day (: `` She told me yes. She related how she once achieved such an intense orgasm that it raised up her kundalini all the way to the top. And the irony, she explained, is that the guy involved knew nothing at all about metaphysics, chakras, or kundalini! ``

\end{sffamily}
\end{footnotesize}

